% Generated from all eight audited P78 conditions; no outcome exclusions.
\section{API失敗通知の通常対照}
同じ４課題を開始から各１回，実API例外を文字列値で返すVとPython例外で通知するEを比較した．
型付き正常値，最初からの推論保持，持続Python，公開名と予算は共通である．
各80生成・450,000 token・1,800秒，APIは写真160・他80とし，過去比較の再現とはしない．
\begin{center}\begin{minipage}{\columnwidth}\centering
\small 失敗通知の全８条件：成否・仕事token\par\smallskip
\begin{tabular}{lrrc}\hline 課題 & V：文字列 & E：例外 & 実例外V/E\\\hline
写真 & 成功 184,501 & 成功 181,555 & 1/1\\
精算 & 成功 214,306 & 成功 213,504 & 3/3\\
返金 & 成功 186,960 & 成功 119,442 & 0/0\\
書出 & 未達 84,946 & 未達 277,203 & 0/1\\
\hline
合計 & 670,713 & 791,704 & 4/5\\
\hline\end{tabular}\end{minipage}\end{center}
両方3/4成功，全使用量既知で，E$-$Vは+120,991 tokenだった．
全件評価・監査・費用既知の下，実例外と双方成功があり，成功を失わずEの合計費用が減る採用条件は未達だった．
初回の実送信要求は全対一致した．参照IDだけをそろえた事後照合では，写真・精算は第３回の例外後，第４回に入力が分かれた．
返金・書出しは第６回に残時間通知だけが１秒異なり，公開出力も分かれた．返金には実例外がなく，書出しEの最初の例外は第19回だった．
ID置換後のJSON照合に限り，通知前の差を通知方法の効果にしない．残時間１秒の因果効果も未特定である．
全137生成・448 API・８評価，推論器全稼働4,390.281秒を記録した．
